\documentclass[10pt,a4paper]{amsart}
\usepackage[margin=19mm]{geometry}
\usepackage{amsmath,amssymb,mathtools}
\usepackage[scr=rsfs,cal=euler]{mathalfa}
\usepackage{xfrac}
\usepackage[unicode,hidelinks,bookmarks=false]{hyperref}
\newcommand{\ie}{i.e.\ }
\newcommand{\cf}{cf.\ }
\newcommand{\resp}{respectively\ }
\newcommand{\Subsection}{\subsection}
\providecommand{\nobreakdash}{\nobreak\hbox{-}}
\setlength{\emergencystretch}{2em}
\multlinegap0pt
\usepackage{amssymb}
\usepackage{tikz}
\newtheorem{theorem}{Theorem}
\newtheorem{proposition}{Proposition}
\title{A Fixed--Point Approach to Non--Commutative Central Limit Theorems}
\author{Jad Hamdan}
\date{Source: arXiv:2305.06960v7 (CC BY 4.0)}
\begin{document}
\maketitle

\noindent\emph{Source and licence.} A Fixed--Point Approach to Non--Commutative Central Limit Theorems. Version arXiv:2305.06960v7 (CC BY 4.0), distributed by arXiv under the Creative Commons Attribution 4.0 International licence, http://creativecommons.org/licenses/by/4.0/, as recorded in the arXiv licence metadata for this exact version and shown on the abstract page https://arxiv.org/abs/2305.06960v7. Full licence notice: \texttt{licenses/arxiv-2305.06960-CC-BY-4.0.txt}.\par\smallskip

\noindent\textit{Editorial scope.} Counted unit: the article's proposition qe:finite1d (source0.tex lines 255--262). The article's complete original proof of it is delivered with it (source0.tex lines 263--278, one \texttt{proof} environment of 16 lines). No mathematics is written, repaired or re-derived: the statement and the proof are the article's own lines, in the article's own notation, and no retained line is substituted or re-flowed.  Retained uncounted as the source-local context the proof needs: lines 146--148; lines 184--189; lines 192--197; lines 215--221.  Nothing was omitted from any retained span, so the package carries no omission record.  No retained line contains a citation, so the wrapper carries no bibliography and no bibliography entry.  No retained line contains a figure environment, an image-inclusion command or drawing code, so no image is delivered and the package declares no asset dependency.  Editorial formatting only, in addition: an AMS \texttt{amsart} preamble that restates the article's own declarations and macros used by the retained lines, namely the environments \texttt{theorem}, \texttt{proposition} on separate counters, together with the standard packages those lines need.  The retained environments print in this document as Theorem 1, Proposition 1 in order of appearance, whereas the article prints the same items with the numbering of its own full document.  The complete native source of the article as distributed in the arXiv e-print for this exact version is bundled unchanged as \texttt{sources/141379/source0.tex}.
\begin{equation}\label{momentcumulantfree}
    m_k(\mu) = \sum_{\substack{\pi \in \mathcal{NC}(k), \\\pi = \{B_1,...,B_n\}} } \prod_{j=1}^n \kappa_{|B_j|}(\mu),\quad \bigg(\text{resp.} \quad   m_k(\mu) = \sum_{\substack{\pi \in \mathcal{I}(k), \\\pi = \{B_1,...,B_n\} }} \prod_{j=1}^n r_{|B_j|}(\mu)\bigg).
\end{equation}

 \begin{theorem}[Benaych-Georges]\label{BGtheorem} Let $p$ be a positive integer and $\mu$ a probability measure on the real line. If $\mu$ admits a $p$--th moment, then $R_\mu$ admits the Taylor expansion
 \[
    R_\mu(z)=\sum_{i=0}^{p-1}\kappa_{i+1}(\mu)z^i+o(z^{p-1})
 \]
 where $(\kappa_n(\mu))_{n\in\mathbb{N}}$ are the free cumulants of $\mu$ \eqref{momentcumulantfree} and the limit is as $z\to 0$ non--tangentially, meaning $|z|\to 0$ and $|\Re(z)|\leq -\alpha \Im{z}$ for some $\alpha>0$.
 \end{theorem}

\begin{proposition}[Arizmendi--Salazar]\label{Kexpansion} Let $p$ be a positive integer and $\mu$ a probability measure on the real line. If $\mu$ admits a $p$--th moment, then $E_\mu$ admits the expansion
\[
    E_\mu(z)=\sum_{i=0}^{p-1}\frac{r_{i+1}(\mu)}{z^i}+o\bigg(\frac{1}{z^{p-1}}\bigg) 
\]
where $(r_n(\mu))_{n\in\mathbb{N}}$ are the Boolean cumulants of $\mu$ and the limit is as $z\to \infty$ non--tangentially.
\end{proposition}

\begin{proposition}\label{Rconv}
        Let $\{\nu_n\}_{n\geq 1}$ be a sequence of Borel probability measures. Then $\nu_n$ converges weakly to a probability measure $\nu$ on $\mathbb{R}$ if and only if 
    \begin{enumerate}
        \item There exists a Stolz angle $\Delta$ such that all $R_{\nu_n}$ are defined on $\Delta$.
        \item $\lim_{n\to\infty}R_{\nu_n}(z)=R_\nu(z)$ for every $z\in \Delta$, and $R_{\nu_n}(-iy)\to 0$ uniformly in $n$ as $y\to 0^+$. 
    \end{enumerate}
\end{proposition}

\begin{proposition}\label{qe:finite1d}
    $d_{\text{\normalfont Free}}$ and $d_{\text{\normalfont Bool}}$ are finite metrics on $\mathcal{M}_0^3$, where convergence in the metric topology implies weak convergence.
    In particular, for any $\mu\in \mathcal{M}_0^3$, there exists constants $B, C>0$ such that
    \begin{align}\label{bound}
         d_{\text{\normalfont Free}}(\mu, \rho_{sc})\leq |m_3(\mu)|+B,\quad d_{\text{\normalfont Bool}}(\mu, \rho_b)\leq |m_3(\mu)|+C.
    \end{align}
    If $\mu\in \mathcal{M}_0^{4}$ and $m_3(\mu)=0$, the claim holds for the $d^{(3)}$ distances as well, replacing $m_3(\mu)$ by $m_4(\mu)$ in the right hand side of the inequalities.
\end{proposition}

\begin{proof} 
    For both $d_{\text{\normalfont Free}}$ and $d_{\text{\normalfont Bool}}$, symmetry is clear, and separation follows from the identity theorem (using complex conjugation to extend to the whole of $\mathbb{C}\setminus \mathbb{R}$) and the fact that probability measures are uniquely determined by their Cauchy/$R$--transform. The triangle inequality follows from $\sup f+g \leq \sup f + \sup g$ and the triangle inequality for the complex norm.

    The fact that convergence in $d_{\text{Free}}$ implies weak convergence is an immediate consequence of theorem \ref{Rconv}. For $d_{\text{Bool}}$, we must also argue that $d_{\text{Bool}}(\mu_n,\mu)\to 0$ implies $E_{\mu_n}\to E_\mu$ on all of $(\mathbb{C}\setminus \mathbb{R})$, for $\mu_n,\mu\in \mathcal{M}_0^3$. To that end, note that for any $R>1$, Chebyshev's inequality gives $\mu_n([-R,R])\geq 1-1/R^2:=C>0$ uniformly over $n$. For any compact $K\subseteq \mathbb{C}\setminus \mathbb{R}$ and $z\in K$, it follows that
    \[
        |\Im G_{\mu_n}(z)|\geq |y|\int_{|t|\leq R} \frac{1}{(x-t)^2+y^2}\mathrm{d}\mu_{n}(t)\geq C \min_{z=x+iy\in K} \frac{|y|}{(|x|+R)^2+y^2} >0
    \]
     uniformly over $n\geq 1$ and $z\in K$. By Montel's theorem $\{E_{\mu_n}\}_n$ is therefore a normal family: on any compact $K\subseteq (C\setminus \mathbb{R})$, every subsequence of $E_{\mu_n}$ has a further converging subsequence to some limit, which must coincide with $E_{\mu}$ by the identity theorem and $d(\mu_n,\mu)\to 0$.

    
    For finiteness of $d_{\text{Free}}$ on $\mathcal{M}_0^3$, note that since the $R$--transform of a probability measure with unit variance is analytic on $|z+i/4|<1/4$, it is bounded on any compact subset of $\Delta_{\epsilon,\epsilon}$ away from zero. It therefore suffices to show that $\lim_{z\to 0} {|R_\mu(z)-R_\nu(z)|}{|z|^{-2}}<\infty$
    for any $\mu,\nu\in\mathcal{M}_0^3$ (with the limit taken inside $\Delta_{\epsilon,\epsilon}$) but this follows immediately from theorem \ref{BGtheorem}. 
    To get (\ref{bound}), we use the same theorem to write $R_\mu(z)=z+m_3(\mu)z^2+z^2v(z)$ for some $v$ satisfying $|v(z)|\to 0$ as $|z|\to 0$, and take $B=\sup_{\Delta_{\epsilon,\epsilon}}|v|<\infty$. Finiteness of $d_{\text{Bool}}$ and the upper bound for $d_{\text{Bool}}(\mu,\rho_b)$ follow by an identical argument, using the expansion for $E_\mu$ (proposition \ref{Kexpansion}).
    
    Lastly, taking one additional term in the expansions of $R_\mu$ and $E_\mu$, the proofs of the claims for $d^{(3)}_{\text{\normalfont{Free}}}$ and $d^{(3)}_{\text{\normalfont{Bool}}}$ are identical.
\end{proof}

\end{document}
